\documentclass[11pt]{article}
\usepackage[margin=1in]{geometry}
\usepackage{amsmath,amssymb,amsthm,mathtools}
\usepackage{array}
\usepackage[hidelinks]{hyperref}

\newtheorem{theorem}{Theorem}[section]
\newtheorem{lemma}[theorem]{Lemma}
\newtheorem{proposition}[theorem]{Proposition}
\newcommand{\spE}{s^+}
\newcommand{\smE}{s^-}
\newcommand{\sig}{\sigma}
\newcommand{\one}{\mathbf 1}

\title{Positive Square Energy of Every Hexacyclic Cactus}
\author{}
\date{26 July 2026}

\begin{document}
\maketitle

\begin{abstract}
Let $G$ be a connected hexacyclic cactus of order $n$.  We prove
$s^+(G)>n$.  Exact block-additive doubly nonnegative optimization leaves only
the cycle multisets $\{3,3,3,3,3,q\}$ and $\{3,3,3,3,5,5\}$.
For disconnected shared-cut graphs we prove a leaf-cluster theorem for the
first family and audit all twenty-eight colored cluster partitions for the
second.  The two entry-sensitive rows are closed by a six-row component
classification and by an exact census of eight incidences and twenty-six
ordered entry orbits.  For one shared-cut cluster, exact color-preserving
incidence censuses contain respectively $68,70,71$ trees in the three
$q$-capacity regimes and $150$ trees in the four-triangle--two-pentagon
family.  All ordinary splits are certified; the bouquets and the unique
saturated pentagon hub are handled directly.  Every construction permits
arbitrary bridge connectors and arbitrary trees attached at arbitrary
vertices.
\end{abstract}

\section{Preliminaries and dependencies}

All graphs are finite, simple, and undirected.  A cactus is a connected graph
whose blocks are bridges or cycles.  Its cyclomatic number is the number of
cyclic blocks.  If the adjacency eigenvalues of $G$ are
$\lambda_1,\ldots,\lambda_n$, put
\[
 \spE(G)=\sum_{\lambda_j>0}\lambda_j^2,
 \qquad \smE(G)=\sum_{\lambda_j<0}\lambda_j^2,
 \qquad \sig(G)=\spE(G)-|V(G)|.
\]
For every vertex partition $V(G)=V_1\mathbin{\dot\cup}\cdots
\mathbin{\dot\cup}V_k$, induced-subgraph superadditivity gives
\begin{equation}\label{eq:super}
 \spE(G)\geq\sum_{i=1}^k\spE(G[V_i]).
\end{equation}
This is the theorem of Akbari, Kumar, Mohar, and Pragada~\cite{AKMP}.

Write $T=C_3$, $P=C_5$, and
\[
 \delta=\sec(\pi/5)-1=\sqrt5-2<\frac12,
 \qquad \delta_q=\sec(\pi/q)-1<1\quad(q\equiv1\pmod4).
\]
We use the following packet bounds.  All allow arbitrary trees attached at
arbitrary vertices, and mixed two-cycle packets allow arbitrary bridge
connectors:
\begin{align}
 \sig(T)&>0,&\sig(P)&\geq-\delta,&
 \sig(C_q)&\geq-\delta_q,\label{eq:one}\\
 \sig(TT)&>1,&\sig(TP)&>1-\delta,&
 \sig(TC_q)&>1-\delta_q.\label{eq:two}
\end{align}
Moreover,
\begin{align}
 \sig(H)&\geq0 &&\text{for every bicyclic or tricyclic cactus }H,
 \label{eq:small}\\
 \sig(H)&>0 &&\text{for every tetracyclic or pentacyclic cactus }H.
 \label{eq:lower-rank}
\end{align}
For the indicated shared-cut clusters we use
\begin{align}
 \sig(TTT)&>2,&\sig(TTTT)&>3,\label{eq:tri-clusters}\\
 \sig(TTP)&>2-\delta
 &&\text{when its two triangles share a cut},\label{eq:ttp}\\
 \sig(TPP)&>6-2\sqrt5,\label{eq:tpp}\\
 \sig(PP)&\geq1-\frac4{3\sqrt{13}}.\label{eq:pp}
\end{align}
Bounds \eqref{eq:one}--\eqref{eq:tpp} and the lower-rank theorems are proved
in the bicyclic, tricyclic, tetracyclic, and pentacyclic companion
manuscripts~\cite{Bicyclic,Tricyclic,Tetracyclic,Pentacyclic};
\eqref{eq:pp} is \cite[Theorem~5.1]{Bicyclic}.
If every cycle is $3\pmod4$ and the cycle-packing number is at most two, the
Sachs half-plane theorem~\cite{PackingTwo} gives, for an $r$-cyclic cactus,
\begin{equation}\label{eq:favorable}
 \sig(H)>r-1.
\end{equation}
This gives the shared $TTT$ estimate.  For four triangles it also applies
unless the incidence is a central triangle with three disjoint triangular
petals.  In that sole case the exact grouped-Sachs matching injection in
\cite[Proposition~4.11]{PackingTwo} gives $s^+>s^-$ and hence
$\sig(TTTT)>3$.  These are the only external packet dependencies.

\section{Connector territories, intervals, and openings}

Join two cyclic blocks in the \emph{shared-cut graph} when they share a cut
vertex, and call its connected components shared-cut clusters.  In the
block-cut tree contract the incidence subtree of each cluster to one marked
node.  Every cluster, including a singleton cyclic block, is marked.  Take the
minimal subtree spanning the marked nodes and suppress only unmarked
degree-two nodes.  The result is the \emph{reduced cluster tree}.

\begin{lemma}[Arbitrary connector territories]\label{lem:territories}
Let $S_1,\ldots,S_k$ be pairwise vertex-disjoint connected subtrees of the
reduced cluster tree, and suppose every marked node belongs to exactly one
$S_i$.  Then $V(G)$ has a partition into connected induced territories
$G_1,\ldots,G_k$ such that $G_i$ retains exactly the cycles assigned to $S_i$.
Cuts occur on actual bridges, and every attached tree is assigned whole to one
territory.
\end{lemma}

\begin{proof}
Assign every remaining reduced-tree vertex to a nearest labelled subtree,
breaking ties by a fixed order.  In a tree each label class is connected.
Every edge between classes expands to a path consisting only of bridge blocks;
cut one actual bridge on it and assign the two remnants to their labels.  At a
Steiner branch, assign the branch vertex and one initial segment to one class
and cut an actual bridge on every other incident segment.  A component outside
the cyclic hull has a unique attachment to it, for two attachments would make
a cycle in the block-cut tree; assign that entire component to the owner of
the attachment.  The resulting sets are connected, induced, disjoint, and
retain exactly the asserted cycles.
\end{proof}

\begin{lemma}[Intervals and private openings]\label{lem:intervals}
Let $C$ be a cycle node of a cycle--cut incidence tree.
\begin{enumerate}
\item If the branches of $I-C$ are assigned to nonempty consecutive intervals
of $C$, one interval per nonempty group of cyclically consecutive marks, these
intervals and their branches form connected induced territories.  Every mark
has one owner and $C$ is retained by none.
\item If $v$ is a private vertex of $C$, put $v$ and all off-cycle branches
rooted there in $F$.  Then $F$ is a nonempty tree with $\sig(F)=-1$.  If
$C-v$ contains all shared cyclic cuts needed by the other cycles, its
complement is a connected induced territory.
\end{enumerate}
Both statements remain true with arbitrary hanging trees.
\end{lemma}

\begin{proof}
Cut the cycle edges between consecutive intervals.  Every interval is a proper
path, and every incidence branch follows its unique mark.  For an opening,
$C-v$ is a path and all branches rooted at $v$ form a tree with it.  Every
nonempty tree $F$ has $s^+(F)=|F|-1$, so $\sig(F)=-1$.  Unique attachment of
hanging trees proves the remaining assertions.  Crossing edges are discarded
only through \eqref{eq:super}; no edge monotonicity is used.
\end{proof}

\section{The exact DNN residual classification}

For a connected edge-containing graph define
\[
 \kappa(G)=\sup_{\substack{M\succeq0,\ M\geq0\\\one^TM\one>0}}
 \frac{4\bigl(\sum_{uv\in E(G)}\sqrt{M_{uv}}\bigr)^2}{\one^TM\one}.
\]
The sharp cactus DNN theorem~\cite{DNN,LTZ} gives, for bridge-block count $b$
and cycle lengths $\ell_1,\ldots,\ell_r$,
\begin{equation}\label{eq:kappa}
 \kappa(G)=b+\sum_j\kappa(C_{\ell_j}),\qquad
 \kappa(C_\ell)=
 \begin{cases}
  \ell,&\ell\text{ even},\\
  \ell\sec^2(\pi/(2\ell)),&\ell\text{ odd},
 \end{cases}
 \qquad \smE(G)\leq\kappa(G).
\end{equation}
For completeness, fold each nonedge entry of a DNN matrix into its two
diagonal entries by adding $M_{uv}(e_u-e_v)(e_u-e_v)^T$.  Rotational averaging
on a cycle and its least Fourier eigenvalue give the displayed cycle constants.
Orthogonal gluing of folded Gram systems at a one-vertex sum, followed by
Cauchy--Schwarz, proves block additivity.  Finally, writing $A=A_+-B$, the
Schur product theorem makes $B\circ B$ DNN and gives
\[
 \smE/2=-\sum_{uv\in E(G)}B_{uv}\leq\sum_{uv\in E(G)}|B_{uv}|,
 \qquad \one^T(B\circ B)\one=\smE;
\]
the definition of $\kappa$ yields $(\smE)^2\leq\kappa(G)\smE$.

Put
\[
 \epsilon_\ell=\begin{cases}0,&\ell\text{ even},\\
 \ell\tan^2(\pi/(2\ell)),&\ell\text{ odd}.
 \end{cases}
\]
Then
\begin{equation}\label{eq:epsilon}
 \epsilon_3=1,
 \quad \epsilon_5=5-2\sqrt5=:a,
 \quad 3a<2,
 \quad 2a>1,
 \quad \epsilon_5+\epsilon_7<1.
\end{equation}
These are exact: $3a<2$ is $169<180$, and $2a>1$ is $80<81$ after
squaring positive sides.  Also
$\cos(\pi/7)>1-\tfrac12(22/49)^2>7/8$, whence
$\epsilon_7<7/15<2\sqrt5-4=1-\epsilon_5$; the middle inequality squares to
$4489<4500$.  The odd sequence decreases because, with
$u=\pi/(2x)$, the derivative of $\tan^2u/u$ is positive exactly when
$2u>\sin u\cos u$.

For a hexacyclic cactus, block counting and \eqref{eq:kappa} give
\[
 b+\sum_{i=1}^6\ell_i=n+5,
 \qquad
 \sig(G)\geq5-\sum_{i=1}^6\epsilon_{\ell_i}.
 \tag{3.1}\label{eq:dnn}
\]
If at most three cycles are triangles, the sum is at most
$3+3a<5$.  With exactly four triangles, an even remaining length contributes
zero; otherwise every pair except $5,5$ contributes at most
$\epsilon_5+\epsilon_7<1$.  The pair $5,5$ survives because $2a>1$.
With at least five triangles the multiset is $\{3,3,3,3,3,q\}$, $q\geq3$.
Thus \eqref{eq:dnn} is strict outside exactly
\begin{equation}\label{eq:residual}
 TTTTTQ=\{3,3,3,3,3,q\}\quad(q\geq3),
 \qquad TTTTPP=\{3,3,3,3,5,5\}.
\end{equation}

\section{A five-triangle shared cluster}

\begin{lemma}[An intersecting-triangle $TTTP$ packet]\label{lem:tttp-margin}
If $T_1,T_2,T_3,P$ form one shared-cut cluster and two of the triangles share
a cut, then $\sig(TTTP)>1$.
\end{lemma}

\begin{proof}
Let $k$ be the number of shared cyclic cuts on $P$.  The incidence excess for
four cycle nodes is three.  The $k$ cuts on $P$ consume at least $k$ units,
and the intersecting triangle pair consumes another unit, either at a cut off
$P$ or as extra degree at a cut on $P$.  Hence $k\leq2$.  Choose a private
vertex $v$ of $P$ and open it as in Lemma~\ref{lem:intervals}.  The path
$P-v$ retains every cut on $P$, so the induced remainder $H$ is connected and
has exactly the three triangular cycles.  The intersecting pair shows that
their cycle-packing number is at most two.  Hence the favorable-cycle theorem
\eqref{eq:favorable}, which does not require one shared-cut cluster, gives
$\sig(H)>2$.  The opened tree has surplus $-1$, and therefore
$\sig(TTTP)>2-1=1$.
\end{proof}

\begin{lemma}\label{lem:five-triangles}
If five triangular blocks form one shared-cut cluster $A$, then
$\sig(A)>2$, including the cycle-packing-three case.
\end{lemma}

\begin{proof}
The bipartite incidence tree has a leaf triangle.  It has one shared cyclic cut,
so open one of its two private vertices by Lemma~\ref{lem:intervals}.  The tree
territory costs one and the remaining four triangles stay connected in one
shared-cut cluster.  Their surplus is greater than three by
\eqref{eq:tri-clusters}; this includes the central-triangle/three-petal
incidence through the exact Sachs matching injection recalled above.  Hence
$\sig(A)>3-1=2$.  In particular no packing-at-most-two assumption is hidden in
the estimate.
\end{proof}

\section{Disconnected shared-cut graph: the $TTTTTQ$ family}

\begin{proposition}\label{prop:disconnected-tttttq}
If the cycles are $TTTTTQ$ and the shared-cut graph is disconnected, then
$\spE(G)>|V(G)|$.
\end{proposition}

\begin{proof}
If $q=3$, the triangular block-graph theorem~\cite{PackingTwo} applies.  Assume
$q\ne3$.  The minimal reduced cluster tree has at least two marked leaves, and
at most one contains $Q$.  Choose a $Q$-free leaf cluster $A$, containing $r$
triangles, and cut the first actual bridge toward the rest.

For $r=1,2$, the remainder is respectively penta- or tetracyclic and is
strictly positive by \eqref{eq:lower-rank}; for $r=3,4$ it is respectively
tri- or bicyclic and is nonnegative by \eqref{eq:small}.  The all-triangle side
is strict, so all these sums are positive.  If $r=5$, Lemma
\ref{lem:five-triangles} gives $\sig(A)>2$.  The remote $Q$ side is
nonnegative unless $q\equiv1\pmod4$, and in that case it is at least
$-\delta_q$ with $\delta_q<1$.  Thus the total is $>2-\delta_q>1$.
Lemma~\ref{lem:territories} makes this proof independent of connector shape and
entry position.
\end{proof}

\section{Disconnected shared-cut graph: all $28$ $TTTTPP$ partitions}

Up to permutations of the four triangles and two pentagons, there are exactly
$29$ colored set partitions of $(4,2)$: recursively choose the first nonzero
pair $(t,p)$ and require the remaining pairs to be lexicographically
nondecreasing.  Removing the one-cluster partition leaves the following
$28$ proper partitions.  The table is both the complete list and the direct
packet audit; the five named rows are proved immediately afterward.
\begin{center}
\renewcommand{\arraystretch}{1.09}
\begin{tabular}{>{\ttfamily}l l}
partition & certificate\\ \hline
T|T|T|T|P|P & Lemma~\ref{lem:all-singletons}\\
TT|T|T|P|P & $>1-2\delta$ plus strict $T$'s\\
TTT|T|P|P & $>2-2\delta$ plus a strict $T$\\
TT|TT|P|P & $>2-2\delta$\\
TTTT|P|P & $>3-2\delta$\\
TP|T|T|T|P & $>1-2\delta$ plus strict $T$'s\\
TTP|T|T|P & Lemma~\ref{lem:ttp-ttp}\\
TT|TP|T|P & $>2-2\delta$ plus a strict $T$\\
TTTP|T|P & Lemma~\ref{lem:e2}\\
TTT|TP|P & $>3-2\delta$\\
TTP|TT|P & $>1-\delta$\\
TTTTP|P & Lemma~\ref{lem:e1}\\
PP|T|T|T|T & shared $PP\geq0$ plus strict $T$'s\\
TT|PP|T|T & $>1$ plus nonnegative $PP$ and strict $T$'s\\
TTT|PP|T & $>2$ plus nonnegative $PP$ and a strict $T$\\
TT|TT|PP & $>2$ plus nonnegative $PP$\\
TTTT|PP & $>3$ plus nonnegative $PP$\\
TPP|T|T|T & $>6-2\sqrt5$ plus strict $T$'s\\
TP|TP|T|T & $>2(1-\delta)$ plus strict $T$'s\\
TTPP|T|T & positive tetracyclic packet plus strict $T$'s\\
TTP|TP|T & $>1-\delta$ plus a strict $T$\\
TPP|TT|T & $>7-2\sqrt5$ plus a strict $T$\\
TTTPP|T & positive pentacyclic packet plus a strict $T$\\
TT|TP|TP & $>3-2\delta$\\
TTTP|TP & $>1-\delta$\\
TTT|TPP & $>8-2\sqrt5$\\
TTPP|TT & positive tetracyclic packet plus $>1$\\
TTP|TTP & Lemma~\ref{lem:double-ttp}
\end{tabular}
\end{center}
Every direct quantity is positive because $\delta<1/2$ and
$6-2\sqrt5>0$.  Whole cluster territories are supplied by
Lemma~\ref{lem:territories}.

\begin{lemma}[All singletons]\label{lem:all-singletons}
The partition $T|T|T|T|P|P$ has positive surplus for every reduced cluster
tree.
\end{lemma}
\begin{proof}
A triangular leaf gives a strict triangle and a positive pentacyclic
remainder.  Otherwise the two pentagons are all the leaves, so the finite tree
has exactly two leaves and is a path.  In path order form $TP+TT+TP$; its
surplus is $>2(1-\delta)+1>0$.
\end{proof}

\begin{lemma}[The double $TTP$ row]\label{lem:double-ttp}
The partition $TTP|TTP$ has positive surplus.
\end{lemma}
\begin{proof}
If the triangles in either cluster meet, \eqref{eq:ttp} plus nonnegativity of
the other tricyclic cluster gives a positive total.  Otherwise each cluster is
the distinct-cut chain $T-P-T$.  Split each central pentagon into two proper
intervals, assigning the bridge connector to the interval containing its
actual entry.  Join those selected intervals through the connector.  The exact
territories are $TT+T+T$, with surplus $>1$ plus two strict terms.  This
coordinated split gives every connector vertex and shared cut one owner.
\end{proof}

\begin{lemma}[The $TTP|T|T|P$ row]\label{lem:ttp-ttp}
Every cactus with cluster partition $TTP|T|T|P$ has positive surplus.
\end{lemma}
\begin{proof}
An intersecting pair in $TTP$ gives $>2-\delta$, which with the remote
pentagon gives $>2-2\delta>0$.  A singleton triangular leaf gives a triangle
and a positive pentacyclic remainder.  Assume neither and write the cluster as
the chain $T-P_0-T$.  If a reduced edge separates $TP$ from a tetracyclic
side, or $TT$ from $TPP$, the corresponding packet sum is positive.  If no
such edge exists, the $TTP$ cluster and remote pentagon are the only leaves;
the two singleton triangles are internal on their path.  Split $P_0$ according
to its actual connector entry and form
\[
 TT+T+TP,
 \qquad \sig>1+0+(1-\delta)>0.
\]
This also covers entries through a triangle, at its attachment cut, or at a
private vertex of $P_0$.
\end{proof}

\begin{lemma}[E1: the six-row $TTTTP|P$ proof]\label{lem:e1}
The partition $TTTTP_0|P_1$ has positive surplus for every connector entry.
\end{lemma}
\begin{proof}
If the four triangles are pairwise disjoint, incidence acyclicity makes $P_0$
a four-petal hub.  An entry through a petal travels with that petal; an entry at
a private hub vertex is paired with an adjacent petal mark.  Splitting $P_0$
gives $TP+TTT$, hence $>1-\delta$.

Suppose two triangles meet.  Delete $P_0$ from the five-cycle incidence tree.
Every resulting component contains triangles, meets $P_0$ at one distinct cut,
and is one component of the triangle shared-cut graph.  At least one has size
at least two, so the component-size partition and the size of the component containing
the connector entry are exactly
\begin{center}
\begin{tabular}{c c c}
sizes & entry size & certificate\\ \hline
$(4)$&$4$&two-pentagon opening\\
$(3,1)$&$3$&$TTTP+T$\\
$(3,1)$&$1$&$TP+TTT$\\
$(2,2)$&$2$&$TTP+TT$\\
$(2,1,1)$&$2$&$TTP+T+T$\\
$(2,1,1)$&$1$&$TP+TT+T$
\end{tabular}
\end{center}
This is exhaustive even for a multiway cut: all triangles at that cut lie in
one component.  For two or three components, cut $P_0$ in one gap between each
consecutive pair of component marks.  Give $P_1$ and its entire connector to
the interval owning the marked entry component.  The five displayed packet
sums are positive by \eqref{eq:one}--\eqref{eq:lower-rank}; in particular
$TTTP+T$ and $TTP+T+T$ are strict although their first term need only be
qualitatively positive or nonnegative.

If instead the connector enters a private vertex $z$ of $P_0$, open $P_0$ at
$z$ and assign the connector and remote pentagon $P_1$ to the rooted territory
at $z$.  This territory is a pentagonal unicyclic cactus and has surplus at
least $-\delta$.  The complementary induced territory contains all four
triangles; the path $P_0-z$ retains every shared cut, and an intersecting pair
ensures that its shared-cut cluster packet sum is greater than one.  The total
is therefore $>1-\delta>0$.

For $(4)$, open private vertices of both pentagons away from the connector
route.  Their tree territories cost two.  The four triangles remain one
shared-cut cluster with surplus $>3$, so the total is $>3-2=1$.  Entries on a
triangle, at an off-$P_0$ triangle cut, at a multiway attachment cut, or through
a rooted branch all belong to one of the six rows and have one owner.
\end{proof}

\begin{lemma}[E2: eight incidences and twenty-six entries]\label{lem:e2}
The partition $TTTP_0|T_4|P_1$ has positive surplus for every reduced-tree
topology and every pair of labelled connector entries.
\end{lemma}
\begin{proof}
If the $T_4$--$P_1$ path avoids the $TTTP_0$ cluster, use $TP$ plus a positive
tetracyclic territory.  Otherwise its two arms enter that cluster.  The exact
colored incidence trees, by cut count $c=1,2,3$, number
\[
 1,\ 3,\ 4.
\]
With triangle nodes $0,1,2$, pentagon node $3$, and cut nodes starting at $4$,
their canonical edge sets are
\begin{align*}
c=1:\quad&((0,4),(1,4),(2,4),(3,4));\\
c=2:\quad&((0,4),(0,5),(1,4),(2,4),(3,5)),\\
&((0,4),(0,5),(1,4),(2,5),(3,4)),\\
&((0,4),(1,4),(2,5),(3,4),(3,5));\\
c=3:\quad&((0,4),(0,5),(0,6),(1,4),(2,5),(3,6)),\\
&((0,4),(0,5),(1,4),(1,6),(2,5),(3,6)),\\
&((0,4),(0,5),(1,4),(2,6),(3,5),(3,6)),\\
&((0,4),(1,5),(2,6),(3,4),(3,5),(3,6)).
\end{align*}
Seven contain a cut incident with two triangles; Lemma~\ref{lem:tttp-margin} gives
$\sig(TTTP_0)>1$, which absorbs $\sig(P_1)\geq-\delta$ independently of both
entries.  The eighth is the pentagon hub with three disjoint triangular petals.

Project an entry through a petal to its hub mark and force it to travel with
that petal; project every other entry to its actual hub vertex.  Thus three
unlabelled distinct petal marks and ordered roots $b$ (for $T_4$) and $p$ (for
$P_1$) lie on five hub vertices, with coincidences allowed.  Modulo the
dihedral group there are exactly $26$ configurations.  Exact enumeration of
all nonempty proper cyclic intervals finds in every configuration an interval
containing $p$ which, after assigning $T_4$ according as $b$ lies in it, owns
one or two triangles.  The certificate count is
\[
 20:\ TP+TTT,
 \qquad 6:\ TTP+TT.
\]
The totals are respectively $>(1-\delta)+0$ and $>0+1$.  Both interval and
complement are nonempty consecutive sets; each petal mark and labelled arm has
one owner.  The executable certificate and its canonical incidences are
specified in Appendix~\ref{app:reproduction}.
\end{proof}

\begin{proposition}\label{prop:disconnected-ttttpp}
If the cycles are $TTTTPP$ and the shared-cut graph is disconnected, then
$\spE(G)>|V(G)|$.
\end{proposition}
\begin{proof}
The twenty-eight-row table is the complete colored cluster partition list.
Twenty-three rows follow directly from the packet ledger; the remaining five
are Lemmas~\ref{lem:all-singletons}--\ref{lem:e2}.  Thus every reduced-tree
topology, connector entry, and colored partition is covered.
\end{proof}

\section{Fully shared $TTTTTQ$: exact census}

For one shared-cut cluster let $I$ be its bipartite cycle--cut incidence tree.
With six cycle nodes and $c$ cut nodes,
\begin{equation}\label{eq:incidence}
 |E(I)|=c+5,
 \qquad \sum_{x\text{ cut node}}(\deg_I(x)-1)=5,
 \qquad 1\leq c\leq5.
\end{equation}
Triangle degrees are at most three; a $q$-cycle degree is at most $q$.

\begin{lemma}[The $TTTTTQ$ census]\label{lem:q-census}
Fix $Q$ and quotient by permutations of the five triangles and cut nodes.  The
numbers of admissible incidence trees are
\begin{center}
\begin{tabular}{c|rrrrr|r}
$Q$ capacity&$c=1$&$c=2$&$c=3$&$c=4$&$c=5$&total\\ \hline
$q=3$&1&6&20&27&14&68\\
$q=4$&1&6&20&28&15&70\\
$q\geq5$&1&6&20&28&16&71
\end{tabular}
\end{center}
Every nonbouquet tree admits a positive ordinary one-cycle interval split:
respectively $67,69,70$ trees.  The sole exception in every regime is
\begin{equation}\label{eq:q-bouquet}
 ((0,6),(1,6),(2,6),(3,6),(4,6),(5,6)),
\end{equation}
where $0,\ldots,4$ are triangles, $5=Q$, and $6$ is the common cut.
\end{lemma}

\begin{proof}
Enumerate cut neighborhoods with $c+5$ incidences, reject zero cycle degree,
cut degree below two, excess capacity, and every non-tree, and canonicalize by
color-preserving permutations.  This gives the displayed counts.  For every
cycle of incidence degree at least two, delete its node and inspect the branch
multisets.  Splitting $Q$ leaves strict all-triangle branches.  Splitting a
triangle leaves one $Q$-branch: $TQ$ is positive; generic $TTQ$ is nonnegative
and another triangle branch is strict; larger $T^kQ$ branches are lower-rank
positive.  If $Q$ is singleton, the other four triangles occupy at most two
branches, so one contains $TT$ and its surplus $>1$ absorbs the hostile loss
$\delta_q<1$.  These tests certify every nonbouquet tree.  The executable
enumeration and asserted canonical exception reproduce the finite step exactly.
\end{proof}

\begin{proposition}\label{prop:shared-tttttq}
If all six $TTTTTQ$ cycles form one shared-cut cluster, then
$\spE(G)>|V(G)|$.
\end{proposition}
\begin{proof}
Realize each split certified by Lemma~\ref{lem:q-census} with consecutive
intervals as in Lemma~\ref{lem:intervals}; the branch packet sum is positive.
For the bouquet, open a private vertex of $Q$ and a private vertex of one
designated triangle.  The path remnants retain the common cut, and the other
four triangles remain one shared-cut cluster.  Thus
\[
 \sig(G)>3-1-1=1.
\]
This also covers $Q=T$: the designation is arbitrary.  The regimes $q=3$,
$q=4$, and $q\geq5$ are exhaustive by \eqref{eq:incidence}.
\end{proof}

\section{Fully shared $TTTTPP$: exact $150$-tree census}

\begin{lemma}[The SAFE incidence census]\label{lem:pp-census}
Modulo permutations of four triangles, two pentagons, and cut nodes, the exact
counts are
\begin{center}
\begin{tabular}{c|rrrrr|r}
$c$&1&2&3&4&5&total\\ \hline
incidence trees&1&9&40&62&38&150\\
SAFE ordinary split&0&9&40&62&37&148
\end{tabular}
\end{center}
The two canonical exceptions, with $0,1,2,3=T$, $4,5=P$, and cuts beginning
at $6$, are
\begin{align}
 &((0,6),(1,6),(2,6),(3,6),(4,6),(5,6));
 \label{eq:pp-bouquet}\\
 &((0,6),(4,6),(1,7),(4,7),(2,8),(4,8),
   (3,9),(4,9),(4,10),(5,10)).
 \label{eq:saturated-hub}
\end{align}
They are respectively the six-cycle bouquet and a saturated pentagon hub with
four triangular petals and one pentagonal petal.
\end{lemma}

\begin{proof}
Equation \eqref{eq:incidence}, cut degree at least two, and cycle capacities
define the finite enumeration.  Canonicalization uses
$S_4\times S_2\times S_c$.  For each of all $6\cdot150=900$ cycle choices,
delete the cycle node and compute the retained incidence components.  The SAFE
ledger is exactly
\begin{center}
\renewcommand{\arraystretch}{1.08}
\begin{tabular}{c l c}
branch&checked retained hypothesis&bound\\ \hline
$T$&none&$>0$\\
$P$&none&$\geq-\delta$\\
$TT$&none&$>1$\\
$TP$&none&$>1-\delta$\\
$PP$&one shared-cut component&$\geq1-4/(3\sqrt{13})$\\
$TTT$&one shared-cut cluster&$>2$\\
$TPP$&one shared-cut cluster&$>6-2\sqrt5$\\
$TTP$&the triangles retain a common cut&$>2-\delta$\\
other rank three&none&$\geq0$\\
$TTTT$&one shared-cut cluster&$>3$\\
$TTTP$&an intersecting triangle pair&$>1$ by Lemma~\ref{lem:tttp-margin}\\
other rank four&none&$>0$\\
rank five&none&$>0$
\end{tabular}
\end{center}
A split is accepted only if the corresponding exact symbolic bound is positive, or equals zero
with a strict summand.  In particular, qualitative lower-rank positivity is
never used to cancel a negative singleton pentagon, and concentrated estimates
are applied only after checking their retained cuts.  The asserted totals are
$1,9,40,62,38$, the accepted totals are $0,9,40,62,37$, and exhaustive
canonicalization leaves exactly \eqref{eq:pp-bouquet} and
\eqref{eq:saturated-hub}.  Appendix~\ref{app:reproduction} gives the executable
certificate.
\end{proof}

\begin{proposition}\label{prop:shared-ttttpp}
If all six $TTTTPP$ cycles form one shared-cut cluster, then
$\spE(G)>|V(G)|$.
\end{proposition}
\begin{proof}
For each of the $148$ SAFE trees, realize the certified split by proper
consecutive intervals and apply its positive packet sum.  In the bouquet,
open private vertices on both pentagons.  Their path remnants retain the common
cut, the four triangles remain one shared-cut bouquet, and
\[
 \sig(G)>3-1-1=1.
\]
In the saturated hub, the hub pentagon uses all five vertices.  In cyclic
order, merge the pentagonal-petal mark with either neighboring mark; both
neighbors are triangular.  Give the remaining three triangular marks separate
proper intervals.  This produces
\[
 TP+T+T+T,
 \qquad \sig(G)>(1-\delta)+0+0+0>0,
\]
with every strict triangle retained.  Thus both exact exceptions are covered.
\end{proof}

\section{Main theorem}

\begin{theorem}\label{thm:main}
Every connected hexacyclic cactus $G$ on $n$ vertices satisfies
\[
 \boxed{\spE(G)>n}.
\]
\end{theorem}

\begin{proof}
The exact DNN estimate \eqref{eq:dnn} is strict unless the cycle multiset is
$TTTTTQ$ or $TTTTPP$.  If the shared-cut graph is disconnected, Propositions
\ref{prop:disconnected-tttttq} and \ref{prop:disconnected-ttttpp} apply.  If it
is connected, all cycles form one incidence-tree cluster and Propositions
\ref{prop:shared-tttttq} and \ref{prop:shared-ttttpp} apply.  These alternatives
exhaust every cycle multiset and every incidence.  Lemmas
\ref{lem:territories} and \ref{lem:intervals} include arbitrary connectors,
Steiner branches, entry vertices, and attached trees.  Every comparison is on
a vertex partition into induced subgraphs; no edge monotonicity or tree
relocation is used.
\end{proof}

Since a connected hexacyclic graph has $|E(G)|=n+5$, this proves the strict
conclusion of the Akbari--Kumar--Mohar--Pragada--Zhang conjecture
\cite[Conjecture~1.2]{AKMPZ} for every hexacyclic cactus.

\appendix
\section{Exact reproduction and canonical proof objects}\label{app:reproduction}

From the repository root run, using only the Python standard library,
\begin{verbatim}
python research/hexacyclic-e2-tttp-entry-census.py
python research/hexacyclic-tttttq-incidence-census.py
python research/hexacyclic-ttttpp-incidence-census.py
\end{verbatim}
The first script asserts the $TTTP$ incidence counts
\verb|{1: 1, 2: 3, 3: 4}|, seven intersecting-pair trees, the unique
three-petal hub, and exactly
\[
 26=20\ (TP+TTT)+6\ (TTP+TT)
\]
ordered labelled-entry orbits.  It prints a canonical incidence, interval, and
packet certificate for each orbit.

The second script asserts
\begin{verbatim}
q=3:  {1: 1, 2: 6, 3: 20, 4: 27, 5: 14}  total 68
q=4:  {1: 1, 2: 6, 3: 20, 4: 28, 5: 15}  total 70
q>=5: {1: 1, 2: 6, 3: 20, 4: 28, 5: 16}  total 71
\end{verbatim}
and the sole canonical exception \eqref{eq:q-bouquet} in each regime.  The
third asserts
\begin{verbatim}
colored trees by cut count: {1: 1, 2: 9, 3: 40, 4: 62, 5: 38}
SAFE one-cycle packet resolutions: {2: 9, 3: 40, 4: 62, 5: 37}
\end{verbatim}
together with all $900$ cycle choices, the exact retained-incidence ledger,
and the two canonical exceptions \eqref{eq:pp-bouquet} and
\eqref{eq:saturated-hub}.  The scripts enumerate finite incidence objects;
the spectral implications are the symbolic packet lemmas proved or cited in
the manuscript.  The last executable uses only exact \texttt{Fraction}
arithmetic.  It replaces radical margins by the rational consequences
$\delta<1/4$, $1-\delta>3/4$, $2-\delta>7/4$,
$6-2\sqrt5>3/2$, and $1-4/(3\sqrt{13})>0$, each certified by squaring positive
rational bounds.  Thus no numerical approximation is a mathematical premise.

The source proof objects integrated into this manuscript are
\begin{verbatim}
research/hexacyclic-dnn-residuals-2026-07-26.md
research/five-triangle-shared-cluster-surplus-2026-07-26.md
research/hexacyclic-tttttq-disconnected-2026-07-26.md
research/hexacyclic-ttttpp-disconnected-audit-2026-07-26.md
research/hexacyclic-ttttpp-e1-resolution-2026-07-26.md
research/hexacyclic-e2-tttp-two-entry-resolution-2026-07-26.md
research/hexacyclic-tttttq-incidence-census.py
research/hexacyclic-ttttpp-incidence-census.py
research/hexacyclic-e2-tttp-entry-census.py
\end{verbatim}

\begin{thebibliography}{99}
\bibitem{AKMP}
S.~Akbari, H.~Kumar, B.~Mohar, and S.~Pragada,
\emph{A linear lower bound for the square energy of graphs},
Electron. J. Combin. 32(3) (2025), Paper No. P3.53.

\bibitem{AKMPZ}
S.~Akbari, H.~Kumar, B.~Mohar, S.~Pragada, and S.~Zhang,
\emph{Refinement of a conjecture on positive square energy of graphs},
arXiv:2506.07264, 2025.

\bibitem{Bicyclic}
Companion manuscript, \emph{Positive Square Energy of Every Bicyclic Cactus},
2026; repository path \texttt{all-bicyclic-cacti/paper.tex}.

\bibitem{Tricyclic}
Companion manuscript, \emph{Positive Square Energy of Every Tricyclic Cactus},
2026; repository path \texttt{all-tricyclic-cacti/paper.tex}.

\bibitem{Tetracyclic}
Companion manuscript, \emph{Positive Square Energy of Every Tetracyclic
Cactus}, 2026; repository path \texttt{all-tetracyclic-cacti/paper.tex}.

\bibitem{Pentacyclic}
Companion manuscript, \emph{Positive Square Energy of Every Pentacyclic
Cactus}, 2026; repository path \texttt{all-pentacyclic-cacti/paper.tex}.

\bibitem{DNN}
Companion manuscript, \emph{The Optimized Liu--Tang--Zhang Constant for Cactus
Graphs}, 2026; repository path \texttt{sharp-cactus-dnn/paper.tex}.

\bibitem{PackingTwo}
Companion manuscript, \emph{Square Energy from Cycle Packing and Block-Graph
Territories}, 2026; repository path
\texttt{packing-two-square-energy/paper.tex}.

\bibitem{LTZ}
Y.~Liu, Q.~Tang, and S.~Zhang,
\emph{The positive and negative square-energy conjecture},
arXiv:2607.18031, 2026.
\end{thebibliography}
\end{document}
